\scanpage{23}
\mspage{\unsure{6}}
\uline{Thm} Considérons les \inserted{deux} catégories \deleted{\unsure{ainsi définies|définies ainsi}} :
a) celle des \deleted{\unsure{couples|triples}}\inserted{systèmes}
% P05: the dotted X has a separate dense upper-right group; no recoverable token
% or cancellation state is established. The empty payload records present ink.
$(\widetilde X^*,\dot X^{\unsure{}},\mathcal E,Y,f)$, où
$\widetilde X^*$ espace top., $\dot X$ partie fermée, $\mathcal E$ un germe
de voisinage cylindrique de $\dot X$ dans $\widetilde X^*$,
$Y$ esp. top, $f:\dot X\mathrel{\relbar\joinrel\relbar}Y$ application continue ;
b) celle des systèmes $(X,Y,\mathcal T)$, où
$X$ esp. top., $Y$ partie fermée, $\mathcal T$ un germe
de voisinage cônique de $Y$ dans $X$. \deleted{\unsure{Comme} Alors}
\deleted{\unsure{le}} Alors le foncteur naturel \inserted{ci-dessus} de a)
dans b) est une équivalence de catégories.

% P11: the transverse main-text strokes are not cancellations of the margin.
% The dense sign before the following word is probably cancelled; its sign
% and state remain local doubts. No separate stop is supplied after that word.
\margin{NB Il y a un énoncé \unsure{analogue} en termes de \unsure{structures} de
voisinages tubulaires \uline{stricts} (au lieu des germes
de voisinages tub.) \unsure{à} \uuline{voir !}
Les catégories sont \unsure{\deleted{$\simeq$}|$\simeq$} \unsure{équivalentes}
Les morphismes ici sont les isos, pour simplifier !}

Rappelons que
\[
X=\widetilde X^*\amalg_{\dot X}Y
\]
% P13: ci-dessus remains doubtful within the marginal note; no independent
% terminal point is visible. P14's main-line insertion leader is kept separate.
\margin{\uline{NB} $(\widetilde X^*,\dot X)$ joue le rôle d'un
``découpage cônique'' de $X$ suivant $(Y,\mathcal T)$ \unsure{ci-dessus}}

\uline{NB} Une partie \inserted{$X'$} de $X$ revient donc à une
partie $\widetilde X'^*$ de $X^*$, $Y'$ de $Y$, ayant $=$
images inverses dans $\dot X$, et $X'$ est fermé ssi
$\widetilde X'^*$ et $Y'$ le sont. \deleted{Pour que $X'$ soit}
\inserted{Pour que $X'$ soit}
compatibles avec le germe de structure
% P18: une description is raised and has its own cancellation stroke.
% The small downward hook after old et belongs to the insertion/leader area;
% no additional letter or abbreviation is recoverable from it.
cônique $\mathcal T$, \deleted{il f. \unsure{et} \inserted{\deleted{une description}} que $X'_{\unsure{*}}$}
\inserted{il f et s. que l'on ait}
a) $\widetilde X'$ compatible avec le germe de structure
cylindrique $\mathcal E$ de $(\widetilde X,\dot X)$
Il y a lieu alors de \unsure{completer|compléter} $\widetilde X'$ par
% P21: this label is raised above-left of the closure formula. No equals
% sign is written between the label and that formula.
$\inserted{\widetilde X'_1}\,\overline{X'\cap X^*}$
(où $X^*=\widetilde X\setminus\dot X$), et le germe de $\widetilde X'_1$ au vois.
\margin{22}
